\section{Relative computation and Post's theorem}
\label{sec:decision-hierarchy-03}

Many-one reducibility compares sets by a total computable translation of
questions.  Relative computation compares them by permitting one
computation to consult a fixed set during its run.

\begin{definition}[Oracle Turing machine and Turing reducibility]
An \emph{oracle Turing machine} with oracle \(B\subseteq\N\) is a
deterministic single-tape Turing machine augmented with a distinguished
query operation: at a query, the machine computes a number \(n\) and
control passes to one of two designated states according as \(n\in B\) or
\(n\notin B\).  The run is otherwise governed by the fixed machine
definition.  An oracle Turing machine \emph{computes} a partial numerical
function by the input and output discipline fixed for ordinary machines,
read verbatim with the oracle present: on the input word
\(\operatorname{bin}(x)\), its value is the numerical output when the run
halts accepting with a valid numerical output, and its value is undefined
when the run rejects, diverges, or accepts without a valid numerical
output.  A set \(A\subseteq\N\) is \emph{Turing reducible} to
\(B\), written \(A\leq_T B\), iff some oracle Turing machine with oracle
\(B\) decides \(A\).  A set is \emph{computably enumerable relative to}
\(B\) iff it is the domain of a partial function computed by some oracle
Turing machine with oracle \(B\).
\end{definition}

Every many-one reduction yields a Turing reduction: translate the input and
query the oracle once.  The converse fails: \(\N\setminus K\leq_T K\) by
a single query with the answer reversed, while
\(\N\setminus K\not\leq_m K\), because computable enumerability is
preserved downward by many-one reductions, \(K\) is computably enumerable,
and \(\N\setminus K\) is not.

\begin{definition}[Turing jump]
\label{def:turing-jump}
Fix an effective numbering of oracle Turing machines analogous to the
fixed machine numbering: the same designated-slot normalization, with the
query rows included in the row list, and with every number that is not a
valid code denoting an oracle machine that diverges on every input.  For
a valid code \(e\) and an oracle \(B\subseteq\N\), let \(\varphi^B_e(x)\)
be the numerical output when the decoded machine with oracle \(B\), on
the input word \(\operatorname{bin}(x)\), halts accepting with a valid
numerical output, and leave \(\varphi^B_e(x)\) undefined when that run
rejects, diverges, or accepts without a valid numerical output.  The
\emph{Turing jump} of \(B\subseteq\N\) is
\[
  B'\defeq\set{e\in\N:\varphi^B_e(e)\text{ is defined}}.
\]
Define \(B^{(0)}\defeq B\) and
\(B^{(n+1)}\defeq(B^{(n)})'\); in particular \(\varnothing^{(n)}\) is the
\(n\)-th jump of the empty set.
\end{definition}

\begin{lemma}[Relativized adequacy]
\label{lem:relativized-adequacy}
For every oracle \(B\subseteq\N\):
\begin{enumerate}[label=(\alph*)]
\item every partial function computed by an oracle Turing machine with
      oracle \(B\) occurs as \(\varphi^B_e\) for some valid code \(e\);
\item there is a fixed oracle machine computing a function
      \(U^B(e,x)\simeq\varphi^B_e(x)\) for all \(e,x\in\N\); and
\item for every \(m,n\geq1\) there is a total computable function
      \(\mathsf{par}_n^m:\N^{m+1}\to\N\), not depending on \(B\), with
      \(\varphi^B_{\mathsf{par}_n^m(e,\bar a)}(\bar x)
      \simeq\varphi^B_e(\bar a,\bar x)\) for every index \(e\), all
      argument tuples, and every oracle \(B\) simultaneously.
\end{enumerate}
\textit{Status: proof sketch.}
\end{lemma}

\begin{proofplan}
Clause (a) relativizes the proof of \cref{thm:enumeration}: renaming
commutes with the step relation with the query rows included, the
surrogate-start and automatic-alphabet arguments are unchanged, and the
designated-slot normalization applies to the extended row list.  Clause
(b) relativizes the universal simulation behind
\cref{thm:universal-computation}, passing each query of the simulated
machine through to the oracle.  Clause (c) relativizes the compiler
construction given for \cref{thm:parameter}, copying the query rows
unchanged; the compiled table transformation does not consult the oracle,
which gives the simultaneous claim.  The omitted verifications are
exactly those of the results being relativized.
\end{proofplan}

By \cref{lem:relativized-adequacy}(a), a set is computably enumerable
relative to \(B\) iff it is the domain of some \(\varphi^B_e\).

\begin{theorem}[Jump theorem]
\label{thm:jump}
For every \(B\subseteq\N\):
\begin{enumerate}[label=(\alph*)]
\item \(B'\) is computably enumerable relative to \(B\); and
\item \(B'\not\leq_T B\).
\end{enumerate}
\textit{Status: proof sketch.}
\end{theorem}

\begin{proofplan}
(a) \(B'\) is the domain of the partial function
\(e\mapsto U^B(e,e)\), which an oracle machine with oracle \(B\) computes
by \cref{lem:relativized-adequacy}(b).

(b) Relativized diagonal.  Suppose some oracle machine with oracle \(B\)
decides \(B'\), and let \(h^B\) be its total characteristic function,
with \(h^B(e)=1\) iff \(e\in B'\).  Define
\[
  d^B(e)\simeq
  \begin{cases}
    0, & h^B(e)=0,\\
    \text{undefined}, & h^B(e)=1.
  \end{cases}
\]
Then \(d^B\) is computed by an oracle machine with oracle \(B\): run the
decider, then write \(0\) and accept, or enter an infinite loop.  By
\cref{lem:relativized-adequacy}(a), fix \(b\) with
\(\varphi^B_b\simeq d^B\).  Then \(b\in B'\) iff \(\varphi^B_b(b)\) is
defined iff \(h^B(b)=0\) iff \(b\notin B'\), a contradiction.

The displayed diagonal calculation is complete, but clause (a) rests on
the relativized universal machine and clause (b) on the occurrence clause
of \cref{lem:relativized-adequacy}, both proof sketches; the status
policy of the front matter forbids complete-proof status over a sketch
dependency, which is why this theorem is a proof sketch.
\end{proofplan}

\begin{theorem}[Post]
For every \(n\in\N\) and every \(A\subseteq\N\):
\begin{enumerate}
  \item \(A\in\boldsymbol{\Sigma}^0_{n+1}\) iff \(A\) is computably
        enumerable relative to \(\varnothing^{(n)}\); and
  \item \(A\in\boldsymbol{\Delta}^0_{n+1}\) iff
        \(A\leq_T\varnothing^{(n)}\).
\end{enumerate}
\textit{Status: cited theorem.}
\end{theorem}

See \cite{Shoenfield1967} for a complete proof; the handbook treatment in
\cite{Barwise1977} develops the relativized machinery in full.  The cited
proofs classify relations through prefix-form formulas;
\cref{thm:prefix-normal-form} transfers those classifications to the
classes fixed here.  Post's
theorem ties the syntactic hierarchy to iterated jumps, and the first
three clauses of the strictness theorem above follow from it together with
the jump theorem, \cref{thm:jump}:
\(\varnothing^{(n+1)}\) is computably enumerable relative to
\(\varnothing^{(n)}\) but not Turing reducible to it, so
\(\boldsymbol{\Sigma}^0_{n+1}\ne\boldsymbol{\Delta}^0_{n+1}\).  The
final inclusion clause additionally uses the join of
\(\varnothing^{(n)}\) with its complement and the completeness of
\(\varnothing^{(n)}\) among sets computably enumerable relative to
\(\varnothing^{(n-1)}\); those steps are part of the cited proof.

\begin{remark}[Convention bridge for the jump]
\label{rem:jump-convention}
Write \(A\equiv_m B\), \emph{many-one equivalence}, iff
\(A\leq_m B\) and \(B\leq_m A\); the standing is a metalinguistic relation
between subsets of \(\N\), and it is used only in this remark and the
results that cite it.  Four computable table transformations align this
jump with the conventions of the cited sources.  (i)
\(K\equiv_m\varnothing'\): an
ordinary transition table embeds as a query-free oracle table, and an
oracle machine with oracle \(\varnothing\) is simulated by answering
every query negatively.  (ii) The halting-based jump of the cited sources
is \(\equiv_m\) the present \(B'\): the map \(e\mapsto e^*\), where
\(e^*\) simulates \(e\) and, on any halt of the simulation, writes \(0\)
and accepts, reduces halting to definedness; the map
\(e\mapsto e^{**}\), where \(e^{**}\) diverges unless the simulated halt
is accepting with a valid numerical output, reduces definedness to
halting.  All four maps are computable table transformations, made
uniform in their inputs by \cref{thm:parameter}.  Jump monotonicity ---
\(A\leq_T B\) implies \(A'\leq_m B'\) --- is used explicitly in this
alignment: replace each \(A\)-query of an oracle machine with oracle
\(A\) by a run of a fixed oracle-\(B\) decider for \(A\), then align the
self-application indices by the relativized parameter theorem,
\cref{lem:relativized-adequacy}(c), with padding.  Hence Post's theorem
is insensitive to the choice between the two conventions, and the
citations stand.
\textit{Status: proof sketch --- the four table transformations are
elementary, but the monotonicity step runs through
\cref{lem:relativized-adequacy}(c), itself a proof sketch.}
\end{remark}

